\section{O efeito da mudança de gestão: estudo de eventos}
\label{sec:efeitogestao}

A forma mais exigente de medir o efeito da gestão OSS é acompanhar
\textbf{os cinco hospitais que de fato mudaram} (Conjunto Hospitalar de
Sorocaba, CNES 2081695; Pérola Byington, CNES 2078287; Hospital Kátia
de Souza Rodrigues de Taipas, CNES 2082225; Hospital Geral de Vila
Penteado, CNES 2091755; e Hospital Estadual de Presidente Prudente,
CNES 2750511) ano a ano em torno da conversão, comparando cada um com
sua própria história e com a evolução simultânea dos demais 284
hospitais. É o chamado \textbf{estudo de eventos}: o ano anterior à
mudança serve de marco zero, e o modelo estima quanto o hospital se
afastou do próprio padrão em cada ano antes e depois da conversão.

As Figuras~\ref{fig:eventomorttmp} e~\ref{fig:eventocustoocup} mostram
os resultados para quatro indicadores. \textbf{Em termos puramente
descritivos}, depois da conversão a
\textbf{mortalidade cai 1,30 ponto percentual já no ano da mudança},
aprofunda para 2,64 e 2,91 pontos nos dois anos seguintes e estabiliza
em 2,23 pontos na janela de três ou mais anos; a leitura causal desses
números depende do teste de tendências anteriores, examinado logo
adiante e rejeitado justamente para a mortalidade. O \textbf{tempo de
permanência} cai de forma crescente até 24,6\% na janela longa. O
\textbf{faturamento real por saída} não reage no ano da mudança, mas acumula
queda de 13,5\% três ou mais anos depois, resultado que deve ser lido
com cautela: o teste de tendências anteriores do faturamento é rejeitado e
40\% dos conversores placebo da Seção~\ref{sec:inferenciafinal}
produzem efeito igual ou maior. \textit{Vale aqui a ressalva do regime
de pagamento enunciada na Seção~\ref{sec:resultados}: para as OSS, o
faturamento por saída é produção sem consequência financeira direta.} A \textbf{ocupação de
internação} sobe após a conversão, com pico no primeiro ano.

Antes de comemorar, é preciso olhar o lado esquerdo das figuras, os
anos \emph{anteriores} à conversão. Ali mora o teste de credibilidade:
se o hospital já vinha melhorando antes de mudar de gestão, parte da
melhora posterior não pode ser atribuída à OSS. O teste formal (os
pontos à esquerda da linha tracejada serem todos estatisticamente
nulos) \textbf{passa para o tempo de permanência (p igual a 0,254) e
para a fração de alta complexidade (p igual a 0,229), é rejeitado para
a mortalidade e a ocupação (p abaixo de 0,001) e é rejeitado também
para o faturamento (p igual a 0,046), ainda que com folga menor}. A rejeição na mortalidade tem nome e
sobrenome, já antecipados na Seção~\ref{sec:conversores}: em Sorocaba
o pico de mortalidade imediatamente anterior à conversão sugere
reversão à média, e no Pérola Byington a melhora começou antes da
mudança formal.

\textbf{O achado que sobrevive a este teste é o tempo de permanência},
e ele merece destaque porque é o único candidato a efeito causal que
esta pesquisa produz. Só ele passa no teste dos anos anteriores com
folga confortável (p igual a 0,254): antes da conversão os hospitais
não vinham encurtando internações, e depois passaram a encurtar. Isso é
o que o desenho de estudo de eventos foi construído para detectar, e é
o que ele detectou. O efeito é de cerca de 10\% de queda, cresce ao
longo dos anos posteriores e reaparece em todas as especificações que
testamos, inclusive na hierárquica da Seção~\ref{sec:fase2} e no
estimador de Callaway e Sant'Anna da
Seção~\ref{sec:inferenciafinal}. Para uma administradora, é a
afirmação mais forte que este trabalho autoriza sobre o efeito de
mudar o modelo de gestão: \textbf{depois da conversão, os mesmos
hospitais passaram a girar seus leitos mais rápido}.

A ressalva vem colada e não pode ser separada dele. Esse efeito é
identificado por \textbf{cinco hospitais}, dois deles com janela
posterior utilizável, e os três testes desenhados para situações de
poucos tratados, permutação, bootstrap selvagem por cluster (uma
técnica de reamostragem que recalcula a margem de erro sorteando
hospitais inteiros milhares de vezes, apropriada quando há poucas
unidades tratadas) e controle sintético,
não o corroboram (Seção~\ref{sec:inferenciafinal}). Também não o
contradizem: com esse número de unidades, eles não distinguem um efeito
nulo de uma ausência de poder para detectá-lo. O achado é, portanto,
\textbf{sugestivo forte, não demonstrado}, e é assim, com as duas
metades juntas, que deve ser citado.

Para a mortalidade e a ocupação, os números medem uma associação
temporal real, mas não um efeito causal limpo.

\comoler{cada ponto é o desvio do hospital convertido em relação ao seu
próprio padrão, no ano indicado; a haste vertical é a margem de erro e
o losango escuro marca o ano de referência, imediatamente anterior à
mudança. À esquerda da linha tracejada o hospital ainda era Direta; à
direita, já era OSS. Efeito crível pede pontos próximos de zero à
esquerda e afastamento só à direita: é o que se vê no tempo de
permanência, mas não na mortalidade, onde os anos anteriores já
oscilam. Log pontos podem ser lidos, aproximadamente, como variação
percentual dividida por 100.}

\begin{figure}[H]
\centering
\subfig{fig_est_01_evento_mort_all.png}{Mortalidade geral, em pontos
percentuais}
\subfig{fig_est_01_evento_tmp.png}{Tempo de permanência, em log pontos
(um log ponto equivale, aproximadamente, a 1\% de variação)}
\caption{Estudo de eventos da conversão para OSS: mortalidade geral e
tempo de permanência, amostra completa.}
\label{fig:eventomorttmp}
\end{figure}

\comofeito{estudo de eventos. O eixo horizontal traz os \textbf{anos
desde a conversão para OSS} (valores negativos antes, positivos depois,
com o ano $-1$ de referência) e o vertical, o \textbf{efeito estimado} do
indicador em relação a esse ano de referência (mortalidade em pontos
percentuais; TMP em log pontos, que equivalem, aproximadamente, à
variação percentual dividida por 100). Cada ponto é o desvio médio dos
conversores naquele ano em relação ao próprio padrão, controlando ano e
hospital; a haste vertical é a margem de erro.}

\begin{figure}[H]
\centering
\subfig{fig_est_01_evento_custo_real.png}{Faturamento real por saída, em log
pontos}
\subfig{fig_est_01_evento_ocupacao_internacao_w.png}{Ocupação de
internação, em log pontos}
\caption{Estudo de eventos da conversão para OSS: faturamento real por saída
e ocupação de internação, amostra completa.}
\label{fig:eventocustoocup}
\end{figure}

\comofeito{mesma construção da figura anterior (anos desde a conversão no
eixo horizontal, efeito estimado no vertical), agora para o faturamento
real por saída e a ocupação de internação, ambos em log pontos.}

\comoler{mesma leitura das figuras anteriores. Nesta figura, nenhum ponto
posterior à conversão se firma longe do zero: a mudança de gestão não
alterou o perfil de complexidade da produção desses hospitais.}

\begin{figure}[H]
\centering
\includegraphics[width=.62\textwidth]{fig_est_01_evento_pct_alta_complex.png}
\caption{Estudo de eventos da conversão para OSS: fração de alta
complexidade, em pontos percentuais, amostra completa.}
\label{fig:eventopctalta}
\end{figure}

\comofeito{mesma construção do estudo de eventos, para a fração de alta
complexidade em pontos percentuais no eixo vertical, ao longo dos anos
desde a conversão no horizontal.}

Como exercício de robustez, o estudo foi refeito na \textbf{região de
suporte comum} (faixas Barcelona 3 e 4, portes médio e grande), onde
Direta e OSS convivem em número suficiente. Sorocaba (faixa 5) e
Presidente Prudente (faixa 2) saem dessa amostra, que fica com Pérola
Byington, Taipas e Vila Penteado. A Figura~\ref{fig:eventosuporte}
mostra o resultado para mortalidade e tempo de permanência: os efeitos
posteriores ficam ainda maiores (mortalidade 2,05 pontos menor no ano
da conversão, tempo de permanência 18,6\% menor), mas o teste dos anos
anteriores é rejeitado em todos os cinco indicadores, e com folga (p
abaixo de 0,001 em todos, com estatística F entre 17,0 e 556,1), e a janela posterior se
apoia em três hospitais, dois deles com um único ano observado após a
mudança. Esses números confirmam a direção dos achados, sem
constituir prova independente.

\comoler{mesma leitura do estudo de eventos anterior, agora só com
hospitais comparáveis entre si. Note que restam menos anos após a
conversão: dos três conversores presentes, dois mudaram em 2025.}

\begin{figure}[H]
\centering
\subfig{fig_est_02_evento_suporte_mort_all.png}{Mortalidade geral, em
pontos percentuais}
\subfig{fig_est_02_evento_suporte_tmp.png}{Tempo de permanência, em log
pontos}
\caption{Estudo de eventos na região de suporte comum (faixas 3 e 4,
portes médio e grande): mortalidade geral e tempo de permanência.}
\label{fig:eventosuporte}
\end{figure}

\comofeito{mesma construção do estudo de eventos (mortalidade e TMP), mas
restrita à \textbf{região de suporte comum} (faixas Barcelona 3 e 4 e
portes médio e grande), onde Direta e OSS convivem em número
suficiente; ficam de fora Sorocaba (faixa 5) e Presidente Prudente
(faixa 2), restando três conversores.}

